\documentclass[12pt, a4paper]{article}

% --- PAQUETES ESENCIALES ---
\usepackage[utf8]{inputenc}
\usepackage[T1]{fontenc}
\usepackage[spanish]{babel}
\usepackage{amsmath}
\usepackage{graphicx}
\usepackage{booktabs}
\usepackage{geometry}
\usepackage{hyperref}
\usepackage{fancyhdr}
\usepackage[table]{xcolor}
\usepackage{listings}

% --- ETIQUETA 'LISTING' A 'LISTADO' ---
\addto{\captionsspanish}{\renewcommand{\lstlistingname}{Listado}}
\addto{\captionsspanish}{\renewcommand{\lstlistlistingname}{Listado de Códigos}}

% --- ESTILO DE CÓDIGO ---
\definecolor{GrisClaro}{rgb}{0.95, 0.95, 0.95}

\lstset{
    backgroundcolor=\color{GrisClaro},
    basicstyle=\ttfamily\small,
    breaklines=true,
    showstringspaces=false,
    numbers=left,
    numberstyle=\tiny\color{gray},
    frame=single,
    framesep=5pt,
    rulesepcolor=\color{black},
    keywordstyle=\color{blue}\bfseries,
    commentstyle=\color{green!60!black}\itshape,
    stringstyle=\color{orange},
    tabsize=2,
}

% Lenguajes que listings no trae nativamente
\lstdefinelanguage{Kotlin}{
    keywords={package, import, fun, override, val, var, when, class, return, super, private, public, internal, object, this, if, else, true, false, null, in, is, as, try, catch},
    keywordstyle=\color{blue}\bfseries,
    ndkeywords={String, Int, Boolean, Unit, Bundle, MethodChannel, FlutterFragmentActivity, FlutterEngine, WindowManager, Settings, Exception},
    ndkeywordstyle=\color{purple}\bfseries,
    sensitive=true,
    comment=[l]{//},
    morecomment=[s]{/*}{*/},
    morestring=[b]",
    morestring=[b]',
}

\lstdefinelanguage{Dart}{
    keywords={import, class, abstract, static, final, const, void, return, if, else, async, await, Future, bool, String, int, double, true, false, null, this, super, extends, new, try, catch, on, in, is, as, override, enum, switch, case, default},
    keywordstyle=\color{blue}\bfseries,
    ndkeywords={Widget, State, StatefulWidget, StatelessWidget, BuildContext, MaterialApp, MethodChannel, Platform, PlatformException, AlertDialog, PopScope, SystemNavigator, showDialog, kDebugMode, UsbDebuggingGuard, UsbDebuggingResult, UsbDebuggingStatus, IntegrityViewModel},
    ndkeywordstyle=\color{purple}\bfseries,
    sensitive=true,
    comment=[l]{//},
    morecomment=[s]{/*}{*/},
    morestring=[b]",
    morestring=[b]',
}

% --- MÁRGENES ---
\geometry{
    a4paper,
    total={170mm,257mm},
    left=25mm,
    right=25mm,
    top=30mm,
    bottom=30mm,
}

% --- ENCABEZADO Y PIE ---
\pagestyle{fancy}
\fancyhf{}
\renewcommand{\headrulewidth}{0pt}
\rfoot{\thepage}
\lfoot{Seguridad de la Información / Miguel Ángel Gutiérrez Gómez}

% --- DATOS DEL DOCUMENTO ---
\newcommand{\nombreuniversidad}{Universidad Politécnica de Chiapas}
\newcommand{\nombrecarrera}{Ingeniería en Software}
\newcommand{\nombremateria}{Seguridad de la Información}
\newcommand{\nombreprofesor}{MC. José Alonso Macías Montoya}
\newcommand{\corte}{Corte 1}
\newcommand{\actividad}{Práctica - Móvil 3 \\ Detección de Depuración USB \\ (Autoprotección en Tiempo de Ejecución / RASP)}
\newcommand{\nombres}{Miguel Ángel Gutiérrez Gómez}
\newcommand{\matricula}{233382}
\newcommand{\fechaentrega}{\today}
\newcommand{\ciudad}{Suchiapa, Chiapas}

\begin{document}

% --- PORTADA ---
\begin{titlepage}
    \centering
    \vspace*{1cm}
    {\Huge\bfseries \nombreuniversidad \par}
    \vspace{0.5cm}
    {\Large \nombrecarrera \par}

    \vspace{2cm}
    \rule{\textwidth}{1.5pt}
    \vspace{0.4cm}
    {\Large\bfseries \actividad \par}
    \vspace{0.4cm}
    \rule{\textwidth}{1.5pt}

    \vspace{3cm}

    \begin{flushleft}
        \centering
        \textbf{Materia:} \nombremateria \\
        \textbf{Profesor:} \nombreprofesor \\
        \textbf{Corte:} \corte
    \end{flushleft}

    \vspace{0.5cm}

    \begin{flushleft}
        \centering
        \textbf{Alumno:} \nombres \\
        \textbf{Matrícula:} \matricula \\
    \end{flushleft}

    \vfill
    \vspace{1cm}
    \textbf{\ciudad, a \fechaentrega}
\end{titlepage}

\clearpage

% =====================================================================
\section{Introducción y Objetivo}
\label{sec:intro}
En el análisis dinámico e ingeniería inversa de aplicaciones móviles, una de
las puertas de entrada más comunes para auditorías no autorizadas es el
\textbf{puente de depuración de Android (ADB)} habilitado mediante la
\textbf{Depuración USB}. Con ADB activo, un atacante puede instrumentar la app,
inspeccionar su memoria, leer logs o conectar herramientas de hooking.

Como medida de \textbf{Autoprotección en Tiempo de Ejecución (RASP)}, esta
práctica integra a la aplicación \textbf{VisionPrice} (Flutter) un módulo que
verifica, durante el arranque, si la Depuración USB está activa. Si lo está y
la app \textit{no} corre en modo desarrollo, se congela la ejecución con un
diálogo mandatorio que instruye al usuario a desactivar la opción y cierra la
aplicación de forma limpia.

% =====================================================================
\section{Arquitectura de la Solución}
\label{sec:arquitectura}
La validación se realiza en el punto más temprano del ciclo de vida con
interfaz: la pantalla \texttt{IntegrityGatePage} (la ruta inicial de la app).
En su \texttt{initState}, mediante \texttt{addPostFrameCallback}, se dispara la
verificación \textbf{antes} de que el usuario pueda alcanzar el Login o
cualquier lógica de negocio. El flujo respeta el patrón MVVM del proyecto:

\begin{itemize}\itemsep2pt
    \item \textbf{Capa de plataforma} — \texttt{UsbDebuggingGuard}
    (\texttt{core/platform/}): contiene la regla de negocio de seguridad,
    aplica la excepción \texttt{kDebugMode} y habla con el canal nativo.
    \item \textbf{Canal nativo} — \texttt{MethodChannel}
    \texttt{visionprice/usb\_debugging} hacia \texttt{MainActivity.kt}, que
    consulta \texttt{Settings.Global.ADB\_ENABLED}.
    \item \textbf{ViewModel} — \texttt{IntegrityViewModel} expone
    \texttt{checkUsbDebugging()} y mantiene la lógica fuera de la vista.
    \item \textbf{Vista} — \texttt{IntegrityGatePage} decide, según el
    resultado, mostrar el \texttt{AlertDialog} persistente o continuar.
\end{itemize}

El siguiente fragmento muestra la consulta nativa en Kotlin (Opción B):

\begin{lstlisting}[language=Kotlin, caption={MainActivity.kt — lectura de ADB\_ENABLED}]
"isUsbDebuggingEnabled" -> {
    try {
        val adbEnabled = Settings.Global.getInt(
            contentResolver,
            Settings.Global.ADB_ENABLED, 0
        ) == 1
        result.success(adbEnabled)
    } catch (e: Exception) {
        result.error("ADB_CHECK_FAILED", e.message, null)
    }
}
\end{lstlisting}

La excepción de entorno de desarrollo y el bloqueo persistente en Dart:

\begin{lstlisting}[language=Dart, caption={Guard con excepción kDebugMode + diálogo no descartable}]
// UsbDebuggingGuard.check()
if (kDebugMode) {                 // produccion: la restriccion SI aplica
  return const UsbDebuggingResult(UsbDebuggingStatus.skippedDebugMode);
}
final enabled = await _channel.invokeMethod<bool>('isUsbDebuggingEnabled');

// IntegrityGatePage — dialogo mandatorio
showDialog(
  context: context,
  barrierDismissible: false,            // no se cierra al tocar fuera
  builder: (ctx) => PopScope(
    canPop: false,                      // bloquea el boton "atras"
    child: AlertDialog(
      title: const Text('Acceso bloqueado por seguridad'),
      content: const Text('Desactiva la Depuracion USB para continuar...'),
      actions: [ ElevatedButton(
        onPressed: () => SystemNavigator.pop(),  // cierre limpio
        child: const Text('Entendido, cerrar aplicacion'),
      )],
    ),
  ),
);
\end{lstlisting}

% =====================================================================
\section{Justificación del Método (Opción B: Código Nativo)}
\label{sec:justificacion}
Se eligió la \textbf{Opción B (MethodChannel hacia Kotlin)} en lugar de un
plugin de la comunidad por las siguientes razones:

\begin{enumerate}\itemsep2pt
    \item \textbf{Fuente de verdad directa.} \texttt{Settings.Global.ADB\_ENABLED}
    es el valor que el propio sistema operativo usa para reflejar el estado de
    la Depuración USB. No se infiere de forma indirecta ni depende de
    heurísticas de terceros.
    \item \textbf{Menor superficie de cadena de suministro.} No se agrega
    ninguna dependencia nueva a \texttt{pubspec.yaml}; menos paquetes externos
    significan menos riesgo de abandono o de código malicioso aguas arriba.
    \item \textbf{Control total} sobre el manejo de errores y sobre el valor
    devuelto, además de reutilizar la infraestructura de \texttt{MethodChannel}
    que la app ya empleaba para \texttt{FLAG\_SECURE}.
\end{enumerate}

% =====================================================================
\section{Repositorio}
\label{sec:repo}
Código fuente completo (app Flutter + microservicios):
\begin{center}
\url{https://github.com/MiguelAngelGutierrezGomez/VisionPrice}
\end{center}

% =====================================================================
\section{Capturas de Funcionamiento}
\label{sec:capturas}

\begin{figure}[h]
    \centering
    \fbox{\parbox[c][5.5cm][c]{0.42\textwidth}{\centering \textit{App ejecutándose con normalidad (Depuración USB desactivada): verificación superada y acceso al Login.}}}
    \caption{Entorno seguro: la app continúa al Login.}
    \label{fig:ok}
\end{figure}

\begin{figure}[h]
    \centering
    \fbox{\parbox[c][5.5cm][c]{0.42\textwidth}{\centering \textit{App bloqueada tras activar la Depuración USB: AlertDialog persistente que cierra la app al confirmar.}}}
    \caption{Entorno inseguro: bloqueo mandatorio por Depuración USB.}
    \label{fig:blocked}
\end{figure}

\noindent\textit{Nota:} para reproducir el bloqueo debe compilarse en modo
release/profile (\texttt{flutter run --release}), ya que en modo debug la
verificación se omite por diseño.

\clearpage

% =====================================================================
\section{Análisis Crítico}
\label{sec:analisis}
\textbf{¿Es este método 100\% infalible contra un atacante avanzado con
herramientas de hooking como Frida?}

\textbf{No.} La técnica eleva la barrera frente a auditorías casuales, pero no
es infalible. Un atacante con Frida (o Xposed/LSPosed) opera \textit{dentro}
del proceso de la app y puede:

\begin{itemize}\itemsep2pt
    \item \textbf{Hookear el resultado.} Interceptar el método nativo o el
    handler del \texttt{MethodChannel} y forzar que siempre devuelva
    \texttt{false}, anulando la detección.
    \item \textbf{Operar sin ADB.} Frida puede inyectarse vía
    \texttt{frida-gadget} embebido o en un dispositivo \textit{rooteado}, sin
    necesidad de tener la Depuración USB activa, por lo que la comprobación ni
    siquiera se dispara.
    \item \textbf{Parchear el binario.} Reescribir el \texttt{smali}/\,DEX para
    saltar la rama de bloqueo.
\end{itemize}

\textbf{Capas adicionales que añadiría} (defensa en profundidad):
\begin{enumerate}\itemsep2pt
    \item \textbf{Detección de root y de Frida}: buscar \texttt{su}, binarios y
    paquetes de Magisk, el puerto/named-pipe por defecto de \texttt{frida-server}
    y bibliotecas \texttt{frida-gadget} cargadas en memoria.
    \item \textbf{Anti-debugging nativo} en C/C++ (NDK):
    \texttt{ptrace(PTRACE\_TRACEME)} y lectura de \texttt{TracerPid} en
    \texttt{/proc/self/status} para detectar depuradores adjuntos.
    \item \textbf{Integridad de la app}: verificación de la firma APK y
    detección de re-empaquetado/parcheo.
    \item \textbf{Validación del lado servidor + ofuscación}: que la lógica
    crítica y las reglas de negocio vivan en el backend (no en el cliente,
    siempre manipulable), \textbf{attestation} con Google Play Integrity API, y
    ofuscación del código (R8/ProGuard) para encarecer la ingeniería inversa.
\end{enumerate}

\noindent\textbf{Conclusión:} ninguna defensa del lado cliente es absoluta. El
valor de RASP no es la imposibilidad, sino \textbf{aumentar el costo y el
tiempo} del atacante; la verdadera robustez surge de combinar varias capas y de
no confiar nunca en el dispositivo como entorno seguro.

\end{document}
